\documentclass[11pt]{article}
% Companion to aqfm_overview.tex (informal). For two columns: add "twocolumn" to the class options.
\usepackage[margin=1in]{geometry}
\usepackage{amsmath,amssymb}
\usepackage{booktabs}
\usepackage{float}
\usepackage{microtype}
\usepackage{xcolor}
\usepackage[authoryear,round]{natbib}
\usepackage[colorlinks=true,linkcolor=blue!50!black,citecolor=blue!50!black,urlcolor=blue!50!black]{hyperref}
\setlength{\parskip}{0.5em}
\setlength{\parindent}{0pt}

\title{\bfseries Does the quantum feature map help on real ophthalmic data?\\[0.3em]
\large The same students, now on five public clinical tables}
\author{Christopher Clark}
\date{Draft, \today}

\begin{document}
\maketitle

% =====================================================================
\section{What we're doing and why}

\textbf{Where we left off.} In the teacher--student experiments (companion overview) we let the quantum map
\emph{make} the labels and asked who could learn them. On those labels there is one region (strong coupling,
a medium anneal, a readout including $\langle X_i\rangle$ and $\langle Y_i\rangle$) where nothing classical,
and none of the map's cheap classical versions, catches up. A follow-up study showed two more things. First,
no test on the data alone can say in advance that a dataset \emph{needs} the map: anything classical that
could detect such signal would itself be a classical model using it. Second, a short ladder of cheaper models
\emph{can} rule datasets out. And on hard labels the map did \emph{not} make, it helped only about 4\% of the
time.

\textbf{The question here.} What happens on real data? We took every public tabular ophthalmic dataset we
could obtain and ran exactly the same students on it. We asked two things:
\begin{enumerate}\setlength{\itemsep}{2pt}
\item \textbf{The ladder.} Does the exact map beat everything cheaper (classical features, nonlinear
learners, and the map's own mean-field and no-coupling versions), on data held back from every choice?
\item \textbf{The feature count.} Even if classical models catch up, do they need \emph{more} features than
the map does to get there? That would be a result even without a win.
\end{enumerate}
Everything below (labels, splits, the rule for choosing qubits, what counts as a win, and our predictions) was
written down before any real dataset was modelled, and committed to the repository first
(\texttt{docs/realdata\_plan.md}).

\textbf{How this differs from the synthetic study.} The map and the students are the same. Four things change:
\begin{itemize}\setlength{\itemsep}{2pt}
\item \textbf{There is no teacher.} We don't know the best achievable AUC, and no student has the ``true
recipe'' on its menu. A student can only be compared with the others.
\item \textbf{Claims are made on held-back data.} Each dataset is split in half once. All choices are made on
one half, and the other half is used only at the end.
\item \textbf{The qubits have to be chosen.} Real tables have more features than we have qubits, so we pick
8 of them.
\item \textbf{The feature-count question is new} (Section~\ref{sec:count}).
\end{itemize}

% =====================================================================
\section{The datasets}

Most public ophthalmic data are images; tabular sets are few. Table~\ref{tab:data} lists the five we used.
Two more were planned but not run: a keratoconus set (no usable file obtained) and an electroretinogram set
(needs signal processing first).

\begin{table}[H]\centering\small
\caption{The five datasets. ``Eyes / patients'': PAPILA and GRAPE have two eyes for most patients, so every
split keeps both eyes of a patient together.}\label{tab:data}
\begin{tabular}{@{}p{2.3cm}p{2.6cm}p{3.3cm}p{6.6cm}@{}}\toprule
dataset & rows & label (positives) & candidate features\\\midrule
DR Debrecen \citep{Antal2014} & 1151 images & signs of diabetic retinopathy (611) & 19 measurements extracted from Messidor fundus images: lesion counts, disc size, macula--disc distance, quality and pre-screening flags\\
PAPILA \citep{Kovalyk2022} & 420 eyes / 210 patients & glaucoma (87 eyes, 47 patients) vs healthy; suspects left out & age, sex, refraction, lens status, IOP, corneal thickness, axial length\\
Myopia \citep{Hosmer2013} & 618 children & myopia within 5 years (81) & age, sex, eye biometry, spherical equivalent, weekly activity hours, parental myopia\\
GRAPE \citep{Huang2023} & 263 eyes / 144 patients & glaucoma progression, PLR2 criterion (40 eyes) & baseline age, sex, IOP, corneal thickness, glaucoma type, 5 RNFL thicknesses, 61 visual-field points\\
FIVES / VascX \citep{Jin2022,VascX2024} & 800 images & AMD, DR or glaucoma (600) vs normal & 43 vessel biomarkers (calibre, density, tortuosity, branching)\\\bottomrule
\end{tabular}
\end{table}

\paragraph{Columns we removed on purpose.}
\begin{itemize}\setlength{\itemsep}{2pt}
\item \textbf{PAPILA:} the visual-field mean defect, which is recorded only for glaucoma patients and so gives
the label away.
\item \textbf{GRAPE:} the number of follow-up visits, which isn't known at baseline (and eyes seen more often
have more chances to be called progressing); also the other two progression definitions.
\item \textbf{Everywhere:} IDs and administrative columns.
\end{itemize}
PAPILA's refraction (sphere, cylinder and the astigmatism \emph{axis}, an angle where $0^\circ$ and
$180^\circ$ are the same) was converted to the standard power vectors $M=S+C/2$,
$J_0=-\tfrac C2\cos2\alpha$, $J_{45}=-\tfrac C2\sin2\alpha$ \citep{Thibos1997}.

% =====================================================================
\section{From a clinical table to 8 qubits}

The same steps were applied to every dataset, in this order.
\begin{enumerate}\setlength{\itemsep}{2pt}
\item \textbf{Clean.} Drop any column more than 20\% missing; turn two-level categories into 0/1.
\item \textbf{Split once, 50/50,} into a \emph{discovery} half and a \emph{confirmation} half, balanced for the
label and keeping each patient on one side.
\item \textbf{Fill gaps} with the discovery half's medians.
\item \textbf{Choose the qubits.} A random forest \citep{Breiman2001} is fitted on the discovery half only, and
its 8 most important features become the 8 qubits. We chose forest importance over single-feature
association because it also credits features that matter only in combination, which is what the map
could exploit.
\item \textbf{Scale} each chosen feature with the discovery half's mean and SD and clip at $\pm3$: the same
input pipeline as the synthetic work, so each feature becomes a field $z_i$ on one qubit.
\end{enumerate}
The map, its settings ($J=c/7$ with $c\in\{0.5,\dots,8\}$, twelve anneal times $\tau$) and the measurements
are exactly those of the companion overview.

% =====================================================================
\section{Students, verdicts and the feature count}

\paragraph{The students} are the same as in the synthetic study: the exact map (540 options), mean-field
(540), no coupling (108), ground state (15), classical feature expansions with the 22 tuned maps, and the
nonlinear learners (boosting, RBF SVM, small neural net), all on the 8 chosen inputs. One student is new and
used for reference only: the nonlinear learners on \emph{all} the dataset's features, which shows how much
we lost by keeping only 8.

\paragraph{Discovery half.} Nested cross-validation (5 outer, 3 inner folds, grouped by patient where needed),
exactly as before. Each dataset gets one verdict, checked in this order:
\begin{enumerate}\setlength{\itemsep}{1pt}
\item \textbf{classical enough:} a classical or nonlinear student is within 0.01 AUC of the exact map, or above;
\item \textbf{mean-field enough:} one of the map's classical versions is within 0.01;
\item \textbf{exact ahead (candidate):} the exact map beats every other student by $\ge0.05$;
\item \textbf{exact slightly ahead:} anything in between (no claim).
\end{enumerate}

\paragraph{The feature count.}\label{sec:count} The map hands the classifier its measured expectation values,
not one number per qubit: with 8 qubits that's 8 (one Pauli type), 24 ($X$, $Y$ and $Z$ on every qubit), 28
(one type of pair) or up to 108. So the fair comparison counts \emph{those}:
\begin{gather}
D_Q=\text{number of map features used},\nonumber\\
D^\ast=\text{fewest classical features whose AUC reaches the map's AUC}-0.01,\nonumber\\
R=D^\ast/D_Q .
\end{gather}
$R>1$ would mean classical needs more features than the map. We also compare at \emph{equal} count,
$G=\mathrm{AUC}_{\text{map}}-\mathrm{AUC}$ of the best classical expansion with at most $D_Q$ features, and
do the same for mean-field, which has exactly the same feature count but is cheap to compute. ``Fewer
features'' is a claim about compactness, never about cost: each quantum feature would need on the order of
1000 measurements on hardware.

\paragraph{Confirmation half.} Every student is tuned on the whole discovery half, refitted, and scored
\emph{once} on the confirmation half; 95\% intervals come from 1000 bootstrap resamples of patients. We
would claim a \textbf{win} only if a discovery candidate also beats everything by $\ge0.05$ on confirmation,
with the interval above zero, and \textbf{compactness} only if $G\ge0.02$ there, with the interval above zero.

% =====================================================================
\section{What we found}

\textbf{Headline.} On none of the five datasets does the exact map beat the cheaper students, and it is not
more compact either: classical expansions match it with as few or fewer features everywhere. The version of
the map \emph{with no coupling at all} does as well as the full one.

\begin{table}[H]\centering\small
\caption{AUC of each student, discovery half (nested cross-validation) $\to$ confirmation half (one held-back
scoring). ``Nonlinear'' uses the 8 qubit features; ``all features'' is the reference.
$^\ast$The verdict our rules give for GRAPE is ``exact slightly ahead'' (0.53 vs 0.51), but nothing in GRAPE
has signal (Section~6).}\label{tab:auc}
\footnotesize\setlength{\tabcolsep}{3.5pt}
\begin{tabular}{@{}lccccccc@{}}\toprule
dataset & exact & mean-field & no coupling & classical & nonlinear & all features & verdict\\\midrule
DR Debrecen & 0.78$\to$0.80 & 0.79$\to$0.80 & 0.79$\to$0.80 & 0.78$\to$0.80 & 0.75$\to$0.74 & 0.79$\to$0.81 & classical enough\\
Myopia      & 0.81$\to$0.90 & 0.81$\to$0.90 & 0.81$\to$0.90 & 0.80$\to$0.89 & 0.74$\to$0.75 & 0.74$\to$0.81 & classical enough\\
FIVES/VascX & 0.81$\to$0.83 & 0.81$\to$0.83 & 0.80$\to$0.83 & 0.81$\to$0.83 & 0.80$\to$0.83 & 0.79$\to$0.77 & classical enough\\
PAPILA      & 0.79$\to$0.57 & 0.85$\to$0.58 & 0.80$\to$0.70 & 0.81$\to$0.74 & 0.80$\to$0.65 & 0.75$\to$0.66 & classical enough\\
GRAPE       & 0.53$\to$0.53 & 0.43$\to$0.50 & 0.43$\to$0.45 & 0.51$\to$0.57 & 0.48$\to$0.52 & 0.46$\to$0.57 & slightly ahead$^\ast$\\\bottomrule
\end{tabular}
\end{table}

\begin{table}[H]\centering\small
\caption{Margin of the exact map over the best other student, and the feature count. Confirmation margins
with the 95\% interval; for $G$, the interval's lower end (compactness would need $G\ge0.02$ with the lower
end above 0).}\label{tab:count}
\footnotesize\setlength{\tabcolsep}{3.5pt}
\begin{tabular}{@{}lccccccc@{}}\toprule
dataset & \shortstack{margin\\(discovery)} & \shortstack{margin\\(confirmation)} & $D_Q$ & $D^\ast$ & $R$ & \shortstack{$G$\\(confirmation)} & \shortstack{lower\\end}\\\midrule
DR Debrecen & $-0.015$ & $-0.003$ [$-0.013$, $0.002$] & 24 & 8 & 0.33 & $+0.008$ & $-0.012$\\
Myopia      & $0.000$ & $-0.001$ [$-0.002$, $0.000$] & 8  & 8 & 1.0  & $+0.001$ & $-0.001$\\
FIVES/VascX & $-0.004$ & $-0.002$ [$-0.031$, $-0.000$] & 8  & 8 & 1.0  & $-0.006$ & $-0.029$\\
PAPILA      & $-0.071$ & $-0.171$ [$-0.312$, $-0.052$] & 24 & 8 & 0.33 & $-0.071$ & $-0.137$\\
GRAPE       & $+0.021$ & $-0.047$ [$-0.232$, $0.109$] & 28 & 8 & 0.29 & \multicolumn{2}{c}{(no signal)}\\\bottomrule
\end{tabular}
\end{table}

\textbf{1. No dataset where the map wins.} On the three datasets with clear signal (DR, myopia, FIVES), every
student lands on the same confirmation AUC, within a few thousandths. On PAPILA the map is clearly behind.

\textbf{2. No compactness.} Classical expansions reach the map's AUC with 8 features, against the map's 8--24,
so $R\le1$ everywhere, and no equal-count gap is anywhere near the 0.02 we would need.

\textbf{3. The coupling contributes nothing here.} The no-coupling student matches the exact map on every
dataset. Its features are a nonlinear transform of each input \emph{separately}, so it is effectively an
additive model. The exact student, accordingly, mostly chose single-qubit readouts. These clinical tables
don't contain the kind of joint structure the coupling (and so the quantum part) would add; an interaction
ladder fitted on the side (additive $\to$ pairwise $\to$ three-way $\to$ unrestricted) also reads three of
them as additive.

\textbf{4. This is what the synthetic work predicted.} The map's advantage appeared only on data it generated
itself, and on hard data it didn't generate it helped about 4\% of the time. On five real clinical tables it
helped zero times. Before running we had predicted no wins (held), at most one compactness result (none) and
``classical enough'' on all five (four of five; the fifth is GRAPE, where nothing has signal).

% =====================================================================
\section{Two things we checked afterwards}

These checks were not planned in advance; we did them because two numbers looked odd.

\textbf{PAPILA: why everything drops on the confirmation half.} Plain logistic regression on the same 8 inputs,
which involves no map at all, shows the same pattern: 0.84 within the discovery half, 0.58 within the
confirmation half, 0.64 trained on discovery and tested on confirmation, but 0.81 the other way round. The
single-feature signal differs between halves (age 0.83 vs 0.68; corneal thickness even reverses direction).
The cause is size: the 87 glaucoma eyes come from only 47 patients, about 24 per half, so the halves can differ
a lot by chance. The splitting worked as intended (no patient is in both halves). The lesson is that PAPILA
is much smaller than its 420 rows suggest.

\textbf{GRAPE: no signal from a single baseline visit.} Logistic regression gets 0.46--0.50 within each half
and across. Cross-validated AUCs below 0.5 are a known artefact when there is nothing to learn and few
positives (about 20 per half). Predicting progression from one baseline visit doesn't seem possible with
these features and this sample.

% =====================================================================
\section{What this does and doesn't show}

\begin{itemize}\setlength{\itemsep}{2pt}
\item \textbf{It shows} that on the public tabular ophthalmic data we could find, the annealed quantum feature
map adds nothing beyond its own uncoupled version, and that this is detectable cheaply, before any quantum
resources are spent: the first rung of the ladder (classical models) already stops every dataset.
\item \textbf{It doesn't show} that the map can never help on clinical data. Five datasets is few, two of them
are small in effect (PAPILA, GRAPE), and all are tabular summaries that tend to be close to additive. Data
with genuinely joint structure (many measurements acting together) are where the map's coupling could
matter, and the ladder is the way to check.
\item \textbf{Choosing 8 qubits loses information.} Forest importance favours continuous over binary features
\citep{Strobl2007}: parental myopia, a known risk factor, was not among myopia's 8 qubits. The all-features
reference shows how much this cost on each dataset.
\item \textbf{FIVES has no patient IDs,} so images from the same person could fall in both halves.
\end{itemize}

% =====================================================================
\bibliographystyle{plainnat}
\begin{thebibliography}{10}\small
\bibitem[Antal and Hajdu(2014)]{Antal2014}
B.~Antal, A.~Hajdu. An ensemble-based system for automatic screening of diabetic retinopathy.
\emph{Knowledge-Based Systems} 60, 2014. \href{https://arxiv.org/abs/1410.8576}{arXiv:1410.8576}.

% TODO: verify author list
\bibitem[Kovalyk et~al.(2022)]{Kovalyk2022}
O.~Kovalyk et~al. PAPILA: Dataset with fundus images and clinical data of both eyes of the same patient for
glaucoma assessment. \emph{Scientific Data} 9, 2022.
\href{https://www.nature.com/articles/s41597-022-01388-1}{link}.

\bibitem[Hosmer et~al.(2013)]{Hosmer2013}
D.~W.~Hosmer, S.~Lemeshow, R.~X.~Sturdivant. \emph{Applied Logistic Regression}, 3rd ed. Wiley, 2013
(myopia data from the Orinda Longitudinal Study of Myopia; R package \texttt{aplore3}).

% TODO: verify author list
\bibitem[Huang et~al.(2023)]{Huang2023}
X.~Huang et~al. GRAPE: A multi-modal dataset of longitudinal follow-up visual field and fundus images for
glaucoma management. \emph{Scientific Data} 10, 2023.
\href{https://www.nature.com/articles/s41597-023-02424-4}{link}.

% TODO: verify author list
\bibitem[Jin et~al.(2022)]{Jin2022}
K.~Jin et~al. FIVES: A fundus image dataset for artificial intelligence based vessel segmentation.
\emph{Scientific Data} 9, 2022. \href{https://www.ncbi.nlm.nih.gov/pmc/articles/PMC9352679/}{link}.

% TODO: add author list
\bibitem[VascX(2024)]{VascX2024}
VascX Models: Model ensembles for retinal vascular analysis from color fundus images.
\href{https://arxiv.org/abs/2409.16016}{arXiv:2409.16016}, 2024.

\bibitem[Thibos et~al.(1997)]{Thibos1997}
L.~N.~Thibos, W.~Wheeler, D.~Horner. Power vectors: an application of Fourier analysis to the description and
statistical analysis of refractive error. \emph{Optometry and Vision Science} 74, 1997.

\bibitem[Breiman(2001)]{Breiman2001}
L.~Breiman. Random forests. \emph{Machine Learning} 45, 2001.

\bibitem[Strobl et~al.(2007)]{Strobl2007}
C.~Strobl, A.-L.~Boulesteix, A.~Zeileis, T.~Hothorn. Bias in random forest variable importance measures:
illustrations, sources and a solution. \emph{BMC Bioinformatics} 8, 2007.
\end{thebibliography}

\end{document}
